% =============================================================================
% Supervised Bayesian Matrix Factorization, YFB parameterization
% Full derivation and update algorithm (eta = (Y F) beta)
% Author:  Andrew Walther
% Date:    November 2026
%
% Unifies the April 2026 per-update derivations in components/ (qL_unsupervised,
% qF_supervised, qBeta, ELBO), the October 15 canonical-scale F update
% (components/canonical_ZF) and the November 5 survival-weight note
% (components/survival_block_weight). Formulas are checked against
% code/fit_ssbmf_yfb.R with canonicalize = TRUE; verify_derivation.R in this
% folder checks the EBNM inputs numerically against the update functions.
% =============================================================================

\documentclass[11pt]{article}
\usepackage[margin=1in]{geometry}
\usepackage{amsmath, amssymb, amsthm}
\usepackage{bm}
\usepackage{xcolor}
\usepackage{hyperref}
\usepackage{booktabs}
\usepackage{array}
\usepackage{tabularx}
\usepackage{tcolorbox}
\tcbuselibrary{skins,breakable}

% ---- Custom commands (same conventions as the LB document) -----------------
\newcommand{\lbar}[2]{\bar{l}_{#1#2}}
\newcommand{\fbar}[2]{\bar{f}_{#1#2}}
\newcommand{\zbar}[2]{\bar{z}_{#1#2}}
\newcommand{\betabar}[1]{\bar{\beta}_{#1}}
\newcommand{\lsec}[2]{\overline{l^2_{#1#2}}}
\newcommand{\fsec}[2]{\overline{f^2_{#1#2}}}
\newcommand{\zsec}[2]{\overline{z^2_{#1#2}}}
\newcommand{\betasec}[1]{\overline{\beta^2_{#1}}}
\newcommand{\E}{\mathbb{E}_q}
\newcommand{\Var}{\mathrm{Var}_q}
\newcommand{\KL}{\mathrm{KL}}
\newcommand{\ws}{w_S}

% ---- Coloured boxes ---------------------------------------------------------
\newtcolorbox{ebnmbox}{breakable,
  colback=blue!5!white, colframe=blue!40!black,
  title=EBNM Problem, fonttitle=\bfseries\small}
\newtcolorbox{derivbox}{breakable,
  colback=gray!6!white, colframe=gray!50!black,
  title=Derivation Steps, fonttitle=\bfseries\small}
\newtcolorbox{keybox}[1][Key Point]{breakable,
  colback=green!5!white, colframe=green!40!black,
  title={#1}, fonttitle=\bfseries\small}
\newtcolorbox{correctionbox}[1][Correction to the April 2026 Derivations]{breakable,
  colback=red!5!white, colframe=red!50!black,
  title={#1}, fonttitle=\bfseries\small}
\newtcolorbox{codebox}[1][Implementation Note]{breakable,
  colback=yellow!5!white, colframe=orange!50!black,
  title={#1}, fonttitle=\bfseries\small}

% ---- Document ---------------------------------------------------------------
\title{\textbf{Supervised Bayesian Matrix Factorization}\\
       \large YFB Parameterization, $\bm{\eta} = (YF)\bm{\beta}$:
       Full Derivation and Update Algorithm}
\author{Andrew Walther}
\date{November 2026}

\begin{document}
\maketitle
\tableofcontents
\newpage

% =============================================================================
\section{Introduction}
% =============================================================================

This document derives the coordinate-ascent variational inference (CAVI)
algorithm for the YFB parameterization of supervised Bayesian matrix
factorization. It follows the same path as the LB document
(\texttt{SSBMF-LB-model-derivations/MF\_Derivations\_UpdateAlgo\_REVISED.pdf}):
the joint model, the objective, the Taylor approximation of the Cox partial
likelihood, one update per parameter, and a one-page algorithm summary
(Section~\ref{sec:algorithm}).

The two parameterizations differ in one place. In LB the Cox linear predictor
is $\eta_i = \bm{l}_i^\top\bm{\beta}$, so the patient scores $L$ carry the
survival information. In YFB the linear predictor is
$\eta_i = \bm{y}_i^\top F\bm{\beta}$, the patient's expression profile projected
onto the gene loadings. The gene loadings $F$ are now the parameter that appears
in both likelihoods, and $L$ enters only the reconstruction.

Every update reduces to an empirical Bayes normal means (EBNM) problem: a
pseudo-observation $x = B/A$ with standard error $s = 1/\sqrt{A}$, where $A$
collects the coefficients of the second moment and $B$ the coefficients of the
first moment in the expected log-likelihood. Solving it with the \texttt{ebnm}
package returns the posterior moments and the fitted prior.

\paragraph{Configuration documented.} The derivation describes
\texttt{fit\_ssbmf\_yfb(canonicalize = TRUE, w\_S = $w_S$)}, the configuration
used for the supervision studies from October 2026 onward (point-Laplace priors
on $L$ and $F$, a normal prior on each $\beta_k$). The earlier benchmark
configuration (\texttt{alpha\_F = 0}) is the special case of
Section~\ref{sec:historical}. Options that were tested but not adopted are in
Appendix~\ref{app:options}. Differences from the April 2026 derivations are
listed in Section~\ref{sec:errata}.

\subsection{Update Procedure}

\textbf{For} $b = 1, 2, \ldots$ until convergence:
\begin{enumerate}
  \item Compute the Cox Taylor working quantities $W_{ii}$ and $z_i$ at the
        current linear predictor.
  \item For $k = 1, \ldots, K$:
  \begin{enumerate}
    \item Update $q(\beta_k)$, $g(\beta_k)$ (survival term only).
    \item Update $q(l_k)$, $g(l_k)$ (genomics term only: the standard EBMF update).
    \item Update $q(f_k)$, $g(f_k)$ (genomics and survival terms: differs from EBMF).
  \end{enumerate}
  \item Update $\tau_j$ (the standard EBMF update).
  \item Rescale each program so its training scores have unit standard
        deviation (Section~\ref{sec:scaling}).
\end{enumerate}

% =============================================================================
\section{Notation and Model Setup}
% =============================================================================

\subsection{Index Conventions}

\begin{center}
\begin{tabular}{cll}
\toprule
Symbol & Dimension & Meaning \\
\midrule
$n,\ p,\ K$ & scalars & patients, genes, programs (latent factors) \\
$i,\ j,\ k$ & & patient, gene and program indices \\
$Y$ & $n\times p$ & expression, column-centered and standardized within platform \\
$L$ & $n\times K$ & patient scores; $l_{ik}$ = score of patient $i$ on program $k$ \\
$F$ & $p\times K$ & gene loadings; $f_{jk}$ = loading of gene $j$ on program $k$ \\
$Z = YF$ & $n\times K$ & program scores; $z_{ik} = \sum_j y_{ij} f_{jk}$ \\
$\bm{\beta}$ & $K\times 1$ & Cox coefficients \\
$\tau$ & $p\times 1$ & gene-specific noise precision \\
$(t_i,\delta_i)$ & & survival time and event indicator \\
$D$ & scalar & number of events, $\sum_i\delta_i$ \\
$\ws$ & scalar & survival weight (Section~\ref{sec:elbo}) \\
\bottomrule
\end{tabular}
\end{center}

\subsection{Posterior Moment Notation}

Bars denote moments under the variational posterior $q$:
\[
  \lbar{i}{k} = \E[l_{ik}], \qquad
  \lsec{i}{k} = \E[l_{ik}^2] = \lbar{i}{k}^2 + \Var(l_{ik}),
\]
and likewise $\fbar{j}{k}$, $\fsec{j}{k}$, $\betabar{k}$, $\betasec{k}$.

\subsection{The Program Scores $Z = YF$}

$Z$ is not a separate parameter. It is a fixed linear function of the observed
expression $Y$ and the loadings $F$, so its distribution under $q$ is induced by
$q(F)$. Under the mean-field posterior the loadings are independent across genes
and programs, which gives
\begin{equation}\label{eq:z-moments}
  \zbar{i}{k} = \sum_j y_{ij}\,\fbar{j}{k}, \qquad
  V^Z_{ik} := \Var(z_{ik}) = \sum_j y_{ij}^2\,\Var(f_{jk}), \qquad
  \zsec{i}{k} = \zbar{i}{k}^2 + V^Z_{ik},
\end{equation}
and $z_{ik}$ and $z_{ik'}$ are independent for $k\ne k'$.

\begin{keybox}[No prior is placed on $Z$]
The prior is on $F$. Because $Z$ is a deterministic function of $F$ given $Y$,
its posterior moments follow from $q(F)$ through Eq.~\eqref{eq:z-moments}. The
question of which prior is induced on $Z$ arises only in a model where $Z$
replaces $L$ as a free parameter, and this model does not do that.
\end{keybox}

\subsection{Models}

\paragraph{Genomics model.}
\begin{equation}\label{eq:genomics}
  Y = LF^\top + E, \qquad E_{ij}\overset{\text{ind}}{\sim}\mathcal{N}(0,\tau_j^{-1}).
\end{equation}

\paragraph{Survival model.}
\begin{equation}\label{eq:survival}
  h_i(t) = h_0(t)\exp(\eta_i), \qquad
  \eta_i = \bm{y}_i^\top F\bm{\beta} = \sum_k z_{ik}\beta_k ,
\end{equation}
with $h_0$ unspecified and the Cox partial likelihood (Breslow handling of ties).
$\beta_k$ is the log hazard ratio per unit of program $k$'s score; after the
scaling of Section~\ref{sec:scaling} that unit is one training standard deviation.

\paragraph{Priors.}
\[
  l_{ik}\overset{\text{iid}}{\sim} g_{l_k}, \qquad
  f_{jk}\overset{\text{iid}}{\sim} g_{f_k}, \qquad
  \beta_k \sim g_{\beta_k},
\]
each estimated by empirical Bayes. In the documented configuration $g_{l_k}$ and
$g_{f_k}$ are point-Laplace (point-exponential gives nonnegative programs) and
$g_{\beta_k}$ is $\mathcal{N}(0,\sigma_k^2)$ with $\sigma_k^2$ estimated.

\paragraph{Joint likelihood.}
\begin{equation}\label{eq:joint}
  p(Y, t, \delta \mid L, F, \bm{\beta}, \tau)
  = \underbrace{p(Y\mid L, F, \tau)}_{\text{genomics}}\;
    \underbrace{p(t,\delta \mid Y, F, \bm{\beta})}_{\text{survival}} .
\end{equation}
The survival model conditions on the observed $Y$, which is a covariate in the
Cox model as well as the outcome of the factorization. In LB the survival model
conditions on the latent $L$ instead.

% =============================================================================
\section{Objective Function}\label{sec:elbo}
% =============================================================================

The objective, maximized over $q = \prod_k q_{l_k}q_{f_k}q_{\beta_k}$, the priors
and $\tau$, is
\begin{equation}\label{eq:elbo}
\begin{split}
  \mathcal{F}(q, g, \tau) ={}&
  \E\!\left[\log p(Y\mid L,F,\tau)\right]
  + \ws\,\E\!\left[\ell_{\text{Cox}}(YF\bm{\beta})\right]\\
  &- \sum_k \KL(q_{l_k}\Vert g_{l_k})
  - \sum_k \KL(q_{f_k}\Vert g_{f_k})
  - \sum_k \KL(q_{\beta_k}\Vert g_{\beta_k}).
\end{split}
\end{equation}

\begin{keybox}[The survival weight $w_S$]
$\ws$ multiplies the Cox log-likelihood. At $\ws = 1$, Eq.~\eqref{eq:elbo} is the
ordinary evidence lower bound and its maximizer approximates the Bayesian
posterior. At $\ws \ne 1$ it is the lower bound for a power (generalized-Bayes)
posterior in which the survival likelihood is raised to the power $\ws$. The
genomics term keeps weight 1. The same $\ws$ appears in every update that
involves survival and in the monitored objective. At $\ws = 0$ survival drops out
of every update, including $\beta$ (Section~\ref{sec:historical}).
\end{keybox}

\paragraph{Genomics term.} With $\bar{R}^2_{ij} := \E[(y_{ij}-\bm{l}_i^\top\bm{f}_j)^2]$,
\begin{gather}
  \E[\log p(Y\mid L,F,\tau)]
  = \sum_j\Bigl[\tfrac{n}{2}\log\tau_j - \tfrac{\tau_j}{2}\sum_i \bar{R}^2_{ij}\Bigr] + C,
  \label{eq:gen-term}\\
  \bar{R}^2_{ij} = \Bigl(y_{ij}-\sum_k\lbar{i}{k}\fbar{j}{k}\Bigr)^2
   + \sum_k\bigl(\lsec{i}{k}\fsec{j}{k}-\lbar{i}{k}^2\fbar{j}{k}^2\bigr). \nonumber
\end{gather}

% =============================================================================
\section{Taylor Approximation of the Cox Partial Likelihood}\label{sec:taylor}
% =============================================================================

The Cox partial log-likelihood $\ell(\bm\eta)$ is not conjugate to any of the
priors. As in the LB document (Section~3, which shows every step), it is
replaced by its second-order expansion around the current linear predictor
$\hat{\bm\eta} = Y\bar F\bar{\bm\beta}$, keeping the diagonal of the Hessian:
\begin{equation}\label{eq:taylor}
  \ell(\bm\eta) \approx \ell(\hat{\bm\eta})
  - \frac12\sum_i W_{ii}\,(z_i - \eta_i)^2 + C,
  \qquad
  z_i = \hat\eta_i + \frac{u_i}{W_{ii}} .
\end{equation}
Here $z_i$ is the working response. With $\hat\theta_i = e^{\hat\eta_i}$, risk-set
sums $\hat S(t) = \sum_{l: t_l\ge t}\hat\theta_l$ and Breslow cumulative hazard
$\hat H(t) = \sum_{\text{event times } s\le t} d(s)/\hat S(s)$, the score and the
weight are
\begin{equation}\label{eq:uW}
  u_i = \delta_i - \hat\theta_i\hat H(t_i), \qquad
  W_{ii} = \hat\theta_i\hat H(t_i).
\end{equation}

\begin{codebox}
The exact diagonal of the negative Hessian is
$\hat\theta_i\hat H(t_i) - \hat\theta_i^2\sum_{s\le t_i} d(s)/\hat S(s)^2$. The
code (\texttt{calc\_cox\_taylor\_yf()}) drops the second term, which is small
relative to the first when risk sets are large, and floors $W_{ii}$ at $10^{-6}$.
Because the scores sum to zero at any $\hat{\bm\eta}$,
$\sum_i W_{ii} = \sum_i \hat\theta_i\hat H(t_i) = D$ before the floor: the total
survival weight equals the number of events. The expansion is computed once per
outer iteration (Step~1 of the algorithm), as in LB.
\end{codebox}

\paragraph{Expected surrogate under $q$.} Taking expectations of Eq.~\eqref{eq:taylor},
\begin{equation}\label{eq:surv-expect}
  \E[\ell(\bm\eta)] \approx \ell(\hat{\bm\eta})
  - \frac12\sum_i W_{ii}\Bigl[(z_i - \bar\eta_i)^2 + \Var(\eta_i)\Bigr] + C,
  \qquad
  \Var(\eta_i) = \sum_k\bigl(\zsec{i}{k}\betasec{k} - \zbar{i}{k}^2\betabar{k}^2\bigr),
\end{equation}
since $z_{ik}\beta_k$ are independent across $k$ under the mean-field posterior.

\paragraph{Partial working response.} For program $k$, remove the other programs'
contributions:
\begin{equation}\label{eq:zmk}
  z^{-k}_i = z_i - \sum_{k'\ne k}\zbar{i}{k'}\,\betabar{k'} .
\end{equation}

% =============================================================================
\section{Update for $q(\beta_k)$: Survival Coefficients}\label{sec:beta}
% =============================================================================

Only the survival term and $\KL(q_{\beta_k}\Vert g_{\beta_k})$ involve $\beta_k$.

\begin{derivbox}
Write $\eta_i = z_{ik}\beta_k + \sum_{k'\ne k} z_{ik'}\beta_{k'}$. Keeping the
terms of $-\tfrac12\E[(z_i-\eta_i)^2]$ that involve $\beta_k$, and using the
independence of $\beta_k$, $z_{ik}$ and the other programs under $q$:
\begin{align*}
  -\tfrac12\E\bigl[(z_i - \eta_i)^2\bigr]
  &= -\tfrac12\Bigl[\E[\beta_k^2]\,\E[z_{ik}^2]
     - 2\,\E[\beta_k]\,\zbar{i}{k}\Bigl(z_i - \sum_{k'\ne k}\zbar{i}{k'}\betabar{k'}\Bigr)\Bigr] + C\\
  &= -\tfrac12\,\betasec{k}\,\zsec{i}{k} + \betabar{k}\,\zbar{i}{k}\,z^{-k}_i + C .
\end{align*}
Multiply by $\ws W_{ii}$ and sum over $i$. The result has the form
$-\tfrac12 A\,\E[\beta_k^2] + B\,\E[\beta_k]$, the expected log-likelihood of a
single normal observation of $\beta_k$.
\end{derivbox}

\begin{equation}\label{eq:AB-beta}
  A_{\beta_k} = \ws\sum_i W_{ii}\,\zsec{i}{k}, \qquad
  B_{\beta_k} = \ws\sum_i W_{ii}\,z^{-k}_i\,\zbar{i}{k}.
\end{equation}

\begin{ebnmbox}
\[
  (\betabar{k},\,\betasec{k},\,\hat g_{\beta_k})
  \;\leftarrow\; \mathrm{EBNM}\!\left(x_k = \frac{B_{\beta_k}}{A_{\beta_k}},\;
  s_k = \frac{1}{\sqrt{A_{\beta_k}}}\right).
\]
$\ws$ cancels in $x_k$ and scales $s_k$ by $\ws^{-1/2}$: a larger weight narrows
the survival likelihood for $\beta_k$ relative to its prior. The term $V^Z_{ik}$
inside $\zsec{i}{k}$ is the error-in-variables correction for uncertainty in the
loadings. The prior $g_{\beta_k} = \mathcal N(0,\sigma_k^2)$ is refit at every
update from this single pseudo-observation, which gives
$\hat\sigma_k^2 = \max(0, x_k^2 - s_k^2)$.
\end{ebnmbox}

% =============================================================================
\section{Update for $q(l_k)$: Patient Scores}\label{sec:L}
% =============================================================================

$L$ does not appear in the Cox model, so only the genomics term and
$\KL(q_{l_k}\Vert g_{l_k})$ involve $l_k$. With the partial residual
\begin{equation}\label{eq:Rk}
  R^{-k}_{ij} = y_{ij} - \sum_{k'\ne k}\lbar{i}{k'}\fbar{j}{k'},
\end{equation}
the terms of Eq.~\eqref{eq:gen-term} that involve $l_{ik}$ are
\[
  -\tfrac12\sum_j\tau_j\Bigl[\E[l_{ik}^2]\,\fsec{j}{k}
     - 2\,\E[l_{ik}]\,R^{-k}_{ij}\fbar{j}{k}\Bigr] + C,
\]
which gives
\begin{equation}\label{eq:AB-L}
  A_{l_k} = \sum_j \tau_j\,\fsec{j}{k}\quad(\text{the same for every patient}), \qquad
  B_{l_{ik}} = \sum_j \tau_j\,R^{-k}_{ij}\,\fbar{j}{k}.
\end{equation}

\begin{ebnmbox}
\[
  (\lbar{\cdot}{k},\,\lsec{\cdot}{k},\,\hat g_{l_k})
  \;\leftarrow\; \mathrm{EBNM}\!\left(x_i = \frac{B_{l_{ik}}}{A_{l_k}},\;
  s = \frac{1}{\sqrt{A_{l_k}}}\right), \qquad i = 1,\ldots,n .
\]
This is the standard EBMF update for scores. Survival reaches $L$ only through
$F$: the current $\fbar{j}{k}$ and $\fsec{j}{k}$ carry whatever survival
information the $F$ update put into them. $g_{l_k}$ is refit at every update.
\end{ebnmbox}

% =============================================================================
\section{Update for $q(f_k)$: Gene Loadings}\label{sec:F}
% =============================================================================

$f_{jk}$ appears in both likelihoods: in the reconstruction through
$\bm l_i^\top\bm f_j$, and in the Cox model through $z_{ik} = \sum_j y_{ij}f_{jk}$.
$F$ is the dual-source parameter of YFB, as $L$ is in LB.

\subsection{Genomics Contribution}

The terms of Eq.~\eqref{eq:gen-term} that involve $f_{jk}$ are
$-\tfrac{\tau_j}{2}\sum_i\bigl[\lsec{i}{k}\,\E[f_{jk}^2] - 2\,R^{-k}_{ij}\lbar{i}{k}\,\E[f_{jk}]\bigr]$,
so
\begin{equation}\label{eq:AB-F-gen}
  A^{\text{gen}}_{f_{jk}} = \tau_j\sum_i\lsec{i}{k}, \qquad
  B^{\text{gen}}_{f_{jk}} = \tau_j\sum_i R^{-k}_{ij}\,\lbar{i}{k}.
\end{equation}

\subsection{Survival Contribution}

\begin{derivbox}
Separate gene $j$ from the other genes of program $k$:
\[
  \eta_i = \beta_k\,y_{ij}f_{jk} + \beta_k\,s_{ij} + \eta^{-k}_i,
  \qquad
  s_{ij} = \sum_{j'\ne j} y_{ij'}f_{j'k},
  \qquad
  \eta^{-k}_i = \sum_{k'\ne k} z_{ik'}\beta_{k'} .
\]
Expand $-\tfrac12\E[(z_i-\eta_i)^2]$ and keep the terms that involve $f_{jk}$.
Under the mean-field posterior $f_{jk}$, $\beta_k$, $s_{ij}$ and $\eta^{-k}_i$ are
independent, so
\begin{align*}
  -\tfrac12\E\bigl[(z_i-\eta_i)^2\bigr]
  &= -\tfrac12\Bigl[\betasec{k}\,y_{ij}^2\,\E[f_{jk}^2]
     + 2\,\betasec{k}\,y_{ij}\,\bar s_{ij}\,\E[f_{jk}]\\
  &\qquad\qquad + 2\,\betabar{k}\,y_{ij}\,\bar\eta^{-k}_i\,\E[f_{jk}]
     - 2\,z_i\,\betabar{k}\,y_{ij}\,\E[f_{jk}]\Bigr] + C\\
  &= -\tfrac12\,\betasec{k}\,y_{ij}^2\,\E[f_{jk}^2]
     + \Bigl[\betabar{k}\,y_{ij}\,z^{-k}_i
     - \betasec{k}\,y_{ij}\,\bar s_{ij}\Bigr]\E[f_{jk}] + C,
\end{align*}
using $z_i - \bar\eta^{-k}_i = z^{-k}_i$ (Eq.~\ref{eq:zmk}) and
$\bar s_{ij} = \zbar{i}{k} - y_{ij}\fbar{j}{k}$. Multiply by $\ws W_{ii}$ and sum
over $i$.
\end{derivbox}

\begin{equation}\label{eq:AB-F-surv}
  A^{\text{surv}}_{f_{jk}} = \betasec{k}\sum_i W_{ii}\,y_{ij}^2, \qquad
  B^{\text{surv}}_{f_{jk}} = \betabar{k}\sum_i W_{ii}\,y_{ij}\,z^{-k}_i
   - \betasec{k}\sum_i W_{ii}\,y_{ij}\bigl(\zbar{i}{k} - y_{ij}\fbar{j}{k}\bigr).
\end{equation}

The first part of $B^{\text{surv}}$ is the signal: the correlation, weighted by
$W$, between gene $j$ and the part of the working response that the other
programs leave unexplained. The second part removes what the other genes of
program $k$ already explain.

\subsection{Combined EBNM Problem}

\begin{equation}\label{eq:AB-F}
  A_{f_{jk}} = A^{\text{gen}}_{f_{jk}} + \ws\,A^{\text{surv}}_{f_{jk}}, \qquad
  B_{f_{jk}} = B^{\text{gen}}_{f_{jk}} + \ws\,B^{\text{surv}}_{f_{jk}}.
\end{equation}

\begin{ebnmbox}
\[
  (\fbar{j}{k},\,\fsec{j}{k}) \;\leftarrow\;
  \text{posterior under } g_{f_k}\text{ for }
  x_j = \frac{B_{f_{jk}}}{A_{f_{jk}}},\; s_j = \frac{1}{\sqrt{A_{f_{jk}}}},
  \qquad j = 1,\ldots,p \text{ in sequence},
\]
then refit $\hat g_{f_k}$ by marginal likelihood from the sweep's $(x_j, s_j)$.
\end{ebnmbox}

\begin{keybox}[Why the update is a sequential sweep over genes]
$B^{\text{surv}}_{f_{jk}}$ depends on $\fbar{j'}{k}$ for $j'\ne j$ through
$\zbar{i}{k}$. Exact coordinate ascent therefore updates one gene at a time and
refreshes $\zbar{\cdot}{k}$ and $V^Z_{\cdot k}$ after each gene:
\[
  \zbar{i}{k} \leftarrow \zbar{i}{k} + y_{ij}\bigl(\fbar{j}{k}^{\text{new}} - \fbar{j}{k}^{\text{old}}\bigr),
  \qquad
  V^Z_{ik} \leftarrow V^Z_{ik} + y_{ij}^2\bigl(\Var(f_{jk})^{\text{new}} - \Var(f_{jk})^{\text{old}}\bigr).
\]
The prior $g_{f_k}$ is held fixed during the sweep, so each gene's posterior is a
closed-form normal-means posterior, and it is refit once the sweep ends. At
$\ws = 0$ the coupling disappears and the sweep equals the usual parallel EBMF
update of $f_k$.
\end{keybox}

\paragraph{Precision-weighted combination.} Because $A$ and $B$ are sums of a
genomics and a survival part, the pseudo-observation is a precision-weighted
average of a genomics-only estimate and a survival-only estimate:
\begin{equation}\label{eq:F-combination}
  x_j = \frac{A^{\text{gen}}\,x^{\text{gen}}_j + \ws A^{\text{surv}}\,x^{\text{surv}}_j}
             {A^{\text{gen}} + \ws A^{\text{surv}}},
  \qquad
  x^{\text{gen}}_j = \frac{B^{\text{gen}}}{A^{\text{gen}}},\quad
  x^{\text{surv}}_j = \frac{B^{\text{surv}}}{A^{\text{surv}}}
\end{equation}
(all quantities for gene $j$ and program $k$). The weight on the survival
estimate is $\rho/(1+\rho)$ with $\rho = \ws A^{\text{surv}}/A^{\text{gen}}$.
Section~\ref{sec:balance} works out the size of $\rho$.

\begin{codebox}
\texttt{fit\_ssbmf\_yfb()} calls \texttt{multimodal\_yfb\_update\_F\_mk()} (or its
C++ version) with working weights $\ws W_{ii}$ and signal
$\ws W_{ii}z^{-k}_i$, so the code computes exactly Eq.~\eqref{eq:AB-F}. A program
with $\max_j A_{f_{jk}} < 10^{-8}$ receives no information from either likelihood.
Its posterior would equal the prior and revive $l_k$ and $\beta_k$ from nothing,
so the code sets $f_k = 0$ with a point-mass prior instead. This prunes the
program while keeping $K$ columns.
\end{codebox}

% =============================================================================
\section{Update for $\tau$: Noise Precision}\label{sec:tau}
% =============================================================================

$\tau$ appears only in the genomics term, since the Cox model uses the observed
$Y$. Setting the derivative of Eq.~\eqref{eq:gen-term} with respect to $\tau_j$
to zero gives the standard EBMF update
\begin{equation}\label{eq:tau}
  \hat\tau_j = \frac{n}{\sum_i\bar{R}^2_{ij}} ,
\end{equation}
with $\bar R^2_{ij}$ from Eq.~\eqref{eq:gen-term} (posterior variances included).

% =============================================================================
\section{Score Standardization}\label{sec:scaling}
% =============================================================================

For any $c_k > 0$, the substitution
\[
  f_{\cdot k}\to f_{\cdot k}/c_k,\qquad
  l_{\cdot k}\to c_k\,l_{\cdot k},\qquad
  \beta_k\to c_k\,\beta_k
\]
leaves $LF^\top$ and $\bm\eta = YF\bm\beta$ unchanged. The reconstruction term
cannot fix $c_k$, but the scale of $z_{\cdot k}$ changes the size of every
survival quantity in the updates. Left free, the scale drifted from one
iteration to the next. This feedback caused the April 2026 instability.

After initialization and at the end of every iteration, each program is
rescaled so that its training scores have unit standard deviation:
\begin{equation}\label{eq:scaling}
  c_k = \mathrm{sd}_i\bigl(\zbar{i}{k}\bigr), \qquad
  \fbar{\cdot}{k}\leftarrow\fbar{\cdot}{k}/c_k,\;
  \fsec{\cdot}{k}\leftarrow\fsec{\cdot}{k}/c_k^2,\;
  \lbar{\cdot}{k}\leftarrow c_k\lbar{\cdot}{k},\;
  \lsec{\cdot}{k}\leftarrow c_k^2\lsec{\cdot}{k},\;
  \betabar{k}\leftarrow c_k\betabar{k},\;
  \betasec{k}\leftarrow c_k^2\betasec{k}.
\end{equation}
The loading prior $g_{f_k}$, which is carried across sweeps, is transformed with
$f$ (a Laplace or exponential rate is multiplied by $c_k$, a normal variance
divided by $c_k^2$). The priors on $L$ and $\beta$ are refit at every update and
are scale families, so they need no transformation. A program whose scores have
collapsed ($c_k < 10^{-6}$) is left unscaled.

\begin{keybox}[What the scaling does and does not change]
It selects one representative from a family of equivalent solutions. It changes
neither the fitted reconstruction, the risk scores, nor the objective. It makes
$\beta_k$ a log hazard ratio per training standard deviation of program $k$'s
score, and it keeps the survival terms on a fixed scale between iterations. The
code argument is \texttt{canonicalize = TRUE}. Earlier notes called the result
the ``canonical $Z_F$''.
\end{keybox}

% =============================================================================
\section{Initialization, Convergence, Pruning and Prediction}\label{sec:other}
% =============================================================================

\paragraph{Initialization.} $\bar L$ and $\bar F$ come from the rank-$K$ SVD of $Y$,
$\bar L = U_K D_K^{1/2}$ and $\bar F = V_K D_K^{1/2}$ (absolute values only when the
prior is nonnegative), with $\overline{L^2} = \bar L^2$ and $\overline{F^2} = \bar F^2$.
Then $\bar{\bm\beta} = \bm 0$, $\hat\tau_j = 1/\widehat{\mathrm{var}}(y_{\cdot j})$, and
the scaling of Section~\ref{sec:scaling} is applied. An optional supervised start
(\texttt{init\_supervised = TRUE}) points program $K$ along the null-model Cox
score direction $Y^\top\bm u_0$.

\paragraph{Convergence.} The monitored objective is Eq.~\eqref{eq:elbo} evaluated
at the end of each iteration: the genomics term Eq.~\eqref{eq:gen-term} at the
updated $\tau$, plus $\ws$ times the survival term Eq.~\eqref{eq:surv-expect}
re-expanded at the current $\bar{\bm\eta}$, minus the KL terms. The fit stops
after iteration $b > 5$ when the relative change in this objective is below
\texttt{tol}. The Taylor surrogate is re-expanded at every iteration and the
priors are refit, so the objective is not guaranteed to increase at every step.

\paragraph{Pruning (optional, \texttt{prune = TRUE}).} After convergence, for each
program compute the change in the monitored objective when it is removed. If the
largest change is nonnegative, remove that program and refit from the remaining
programs' moments.

\paragraph{Prediction.} For a new patient with expression $\bm y^*$, processed with
the training transformations,
\begin{equation}\label{eq:predict}
  \hat\eta^* = \bm y^{*\top}\bar F\,\bar{\bm\beta} = \sum_k \bar z^*_k\,\betabar{k},
  \qquad \bar z^*_k = \bm y^{*\top}\fbar{\cdot}{k}.
\end{equation}
The training and test formulas are identical. No refit or posterior for the new
patient's scores is needed (\texttt{predict\_ssbmf\_yfb()}).

% =============================================================================
\section{Update Schedule: What Is Held Fixed, Rescaled, and Updated Together}\label{sec:schedule}
% =============================================================================

Each update maximizes $\mathcal F$ over one block with everything else held at
its current value. The table records the block, when it is computed, and what
it treats as fixed. Any order of exact coordinate updates is valid coordinate
ascent. The order affects the path and possibly the stationary point reached,
not the update formulas.

\begin{center}
\begin{tabularx}{\textwidth}{>{\raggedright\arraybackslash}p{3.0cm}>{\raggedright\arraybackslash}p{3.6cm}>{\raggedright\arraybackslash}X}
\toprule
Quantity & When computed & Held fixed / how updated \\
\midrule
$W_{ii}$, $z_i$ (Taylor) & Once, at the start of each outer iteration, at $\hat{\bm\eta} = Y\bar F\bar{\bm\beta}$ & Frozen for the whole $k$-loop, even after $\bar F$ and $\bar{\bm\beta}$ change. Re-expanded at the end of the iteration for the objective only \\
$\bar z_{\cdot k'}$, $V^Z_{\cdot k}$ & At the start of each program $k$ & Reflect the programs already updated in this sweep (Gauss--Seidel over programs) \\
$z^{-k}_i$ & At the start of each program $k$ & Uses current $\bar z_{\cdot k'}$ and $\betabar{k'}$ of the other programs \\
$q(\beta_k)$ & One program at a time (scalar EBNM) & Uses the frozen $W, z$ and current $\bar Z$. Prior refit at every update (per-program mode) \\
$q(l_k)$ & All $n$ patients together (one EBNM call) & Exact: given $F$ and $\tau$, the patients' scores are conditionally independent and share one precision. Prior refit at every update \\
$q(f_k)$ & Genes one at a time, in order & Exact coordinate ascent: each gene sees the refreshed $\bar z_{\cdot k}$, $V^Z_{\cdot k}$. $g_{f_k}$ fixed during the sweep, refit after it \\
Order within $k$ & $\beta_k$, then $l_k$, then $f_k$ & $f_k$ uses the new $\betabar{k}$, $\betasec{k}$ and $\bar l_{\cdot k}$. (LB uses $l_k$, $f_k$, $\beta_k$) \\
$\tau$ & After the $k$-loop, all genes together & Closed form; genes are independent given $L$, $F$ \\
Objective & After $\tau$ & Computed before rescaling. The rescaling leaves it unchanged \\
Score standardization & End of every iteration (and after initialization) & All programs at once; moments and $g_{f_k}$ transformed (Section~\ref{sec:scaling}) \\
Stopping & After rescaling & Only after iteration 5 \\
\bottomrule
\end{tabularx}
\end{center}

\begin{keybox}[Two approximations in the schedule]
\begin{enumerate}
  \item The Taylor point is not refreshed inside the $k$-loop. Within an
        iteration the updates optimize a fixed quadratic surrogate, and the
        surrogate moves between iterations. This is the usual
        iteratively-reweighted scheme and the reason the objective can dip.
  \item The prior on $f_k$ is fixed during the gene sweep. Refitting it after
        each gene would make the sweep depend on gene order. Holding it fixed
        keeps each gene's update a closed-form posterior.
\end{enumerate}
\end{keybox}

\newpage
% =============================================================================
\section{Complete Algorithm Summary}\label{sec:algorithm}
% =============================================================================

\begin{tcolorbox}[colback=gray!5, colframe=black!60,
                  title=Supervised Matrix Factorization (YFB) --- CAVI Algorithm,
                  fonttitle=\bfseries]

\textbf{Inputs:} $Y_{n\times p}$ (centered, standardized), $(t_i,\delta_i)_{i=1}^n$,
number of programs $K$, survival weight $\ws$ (default 1)

\textbf{Initialize:} $\bar L,\bar F$ from the rank-$K$ SVD of $Y$; $\bar{\bm\beta}=\bm 0$;
$\hat\tau_j = 1/\widehat{\mathrm{var}}(y_{\cdot j})$; rescale so $\mathrm{sd}(Y\fbar{\cdot}{k}) = 1$

\medskip
\textbf{For} $b = 1, 2, \ldots$ \textbf{until convergence:}

\smallskip
\quad\textbf{1.\ Survival working quantities} (once per outer iteration):
\[
  \hat{\bm\eta} = Y\bar F\bar{\bm\beta},\quad
  (u_i, W_{ii}) = \text{CoxTaylor}(\hat{\bm\eta}, \bm t, \bm\delta),\quad
  z_i = \hat\eta_i + u_i/W_{ii}
\]

\quad\textbf{2.\ For} $k = 1,\ldots,K$:
\[
  \zbar{i}{k'} = \textstyle\sum_j y_{ij}\fbar{j}{k'},\quad
  V^Z_{ik} = \textstyle\sum_j y_{ij}^2\,\Var(f_{jk}),\quad
  z^{-k}_i = z_i - \textstyle\sum_{k'\ne k}\zbar{i}{k'}\betabar{k'}
\]

\qquad\textbf{(a)} $q_{\beta_k}$, $g_{\beta_k}$ (survival only):
\[
  A = \ws\textstyle\sum_i W_{ii}(\zbar{i}{k}^2 + V^Z_{ik}),\quad
  B = \ws\textstyle\sum_i W_{ii}\,z^{-k}_i\,\zbar{i}{k},\quad
  (\betabar{k},\betasec{k}) \leftarrow \mathrm{EBNM}(B/A,\,1/\sqrt A)
\]

\qquad\textbf{(b)} $q_{l_k}$, $g_{l_k}$ (genomics only), with
$R^{-k}_{ij} = y_{ij} - \sum_{k'\ne k}\lbar{i}{k'}\fbar{j}{k'}$:
\[
  A = \textstyle\sum_j\tau_j\fsec{j}{k},\quad
  B_i = \textstyle\sum_j\tau_j R^{-k}_{ij}\fbar{j}{k},\quad
  (\lbar{i}{k},\lsec{i}{k}) \leftarrow \mathrm{EBNM}(B_i/A,\,1/\sqrt A)
\]

\qquad\textbf{(c)} $q_{f_k}$ (genomics $+$ survival), genes $j = 1,\ldots,p$ in
sequence with $g_{f_k}$ fixed:
\begin{align*}
  A_j &= \tau_j\textstyle\sum_i\lsec{i}{k} \;+\; \ws\,\betasec{k}\textstyle\sum_i W_{ii}y_{ij}^2\\
  B_j &= \tau_j\textstyle\sum_i R^{-k}_{ij}\lbar{i}{k}
     \;+\; \ws\Bigl[\betabar{k}\textstyle\sum_i W_{ii}y_{ij}z^{-k}_i
     - \betasec{k}\textstyle\sum_i W_{ii}y_{ij}(\zbar{i}{k} - y_{ij}\fbar{j}{k})\Bigr]\\
  &(\fbar{j}{k},\fsec{j}{k}) \leftarrow \text{posterior under } g_{f_k}
    \text{ for } (B_j/A_j,\,1/\sqrt{A_j});\quad
    \text{update } \zbar{\cdot}{k},\,V^Z_{\cdot k}
\end{align*}
\qquad\phantom{\textbf{(c)}} then refit $g_{f_k}$ from $\{(B_j/A_j,\,1/\sqrt{A_j})\}_{j=1}^p$

\smallskip
\quad\textbf{3.\ Noise precision:}
\[
  \bar R^2_{ij} = \Bigl(y_{ij} - \textstyle\sum_k\lbar{i}{k}\fbar{j}{k}\Bigr)^2
  + \textstyle\sum_k\bigl(\lsec{i}{k}\fsec{j}{k} - \lbar{i}{k}^2\fbar{j}{k}^2\bigr),
  \qquad \hat\tau_j = n\big/\textstyle\sum_i\bar R^2_{ij}
\]

\quad\textbf{4.\ Objective:} $\mathcal F$ = genomics term $+\ \ws\,\times$ survival
term (re-expanded at $Y\bar F\bar{\bm\beta}$, with $\Var(\eta_i)$) $-$ KL terms

\quad\textbf{5.\ Score standardization:} $c_k = \mathrm{sd}_i(Y\fbar{\cdot}{k})$;
$\bar f_{\cdot k}/c_k$, $\bar l_{\cdot k}\,c_k$, $\betabar{k}\,c_k$ (second moments by $c_k^{\mp2}$); transform $g_{f_k}$

\quad\textbf{6.\ Convergence:} stop when $b > 5$ and
$|\mathcal F^{(b)} - \mathcal F^{(b-1)}|/|\mathcal F^{(b-1)}| < \text{tol}$

\medskip
\textbf{Prediction:} $\hat\eta^* = \bm y^{*\top}\bar F\bar{\bm\beta}$
\end{tcolorbox}

\paragraph{Where survival enters.}
\begin{center}
\begin{tabular}{lccl}
\toprule
Update & Genomics & Survival & EBNM input \\
\midrule
$q(\beta_k)$ & -- & $\ws\times$ & scalar $(x_k, s_k)$ \\
$q(l_k)$ & \checkmark & -- & $n$-vector, common $s$ \\
$q(f_k)$ & \checkmark & $\ws\times$ & $p$-vector, sequential over genes \\
$\tau$ & \checkmark & -- & closed form \\
\bottomrule
\end{tabular}
\end{center}

% =============================================================================
\section{Balance of the Genomics and Survival Terms}\label{sec:balance}
% =============================================================================

Survival influences the factorization only through the $F$ update, and there
only through the share of the precision it contributes. By
Eq.~\eqref{eq:F-combination} that share is $\rho_{jk}/(1+\rho_{jk})$, with
\begin{equation}\label{eq:rho}
  \rho_{jk} = \frac{\ws A^{\text{surv}}_{f_{jk}}}{A^{\text{gen}}_{f_{jk}}}
  = \frac{\ws\,\betasec{k}\sum_i W_{ii}\,y_{ij}^2}{\tau_j\sum_i\lsec{i}{k}} .
\end{equation}
Eq.~\eqref{eq:rho} is exact and is computed in the code
(\texttt{trace\_supervision = TRUE}). The ratio is invariant to the scale
substitution of Section~\ref{sec:scaling}, since numerator and denominator both
scale by $c_k^2$. Its size can be read off approximately.

\paragraph{Numerator.} With standardized expression and $W_{ii}$ not strongly
related to $y_{ij}^2$, $\sum_i W_{ii}y_{ij}^2 \approx \sum_i W_{ii} = D$
(Section~\ref{sec:taylor}). The numerator is about $\ws\,\betasec{k}\,D$: it
depends on the number of events and on $\beta_k$ per standard deviation, and not
on $p$.

\paragraph{Denominator.} Let $d_k^2 = \lVert\bar l_{\cdot k}\rVert^2\lVert\bar f_{\cdot k}\rVert^2$
be the energy of program $k$'s rank-one component and $v_k = d_k^2/(np)$ its
share of the total variance of $Y$. When the programs are close to orthogonal
and program $k$'s scores are dominated by signal,
$\bar z_{\cdot k} = Y\bar f_{\cdot k} \approx \bar l_{\cdot k}\lVert\bar f_{\cdot k}\rVert^2$.
Unit standard deviation then means
$\lVert\bar l_{\cdot k}\rVert\,\lVert\bar f_{\cdot k}\rVert^2 \approx \sqrt n$,
and combining with the definition of $d_k$ gives
$\lVert\bar l_{\cdot k}\rVert^2 \approx d_k^4/n = n\,p^2 v_k^2$. Hence
\begin{equation}\label{eq:rho-approx}
  \rho_{jk} \;\approx\; \frac{\ws\,\betasec{k}\,(D/n)}{\tau_j\,p^2\,v_k^2}.
\end{equation}

\begin{keybox}[What Eq.~\eqref{eq:rho-approx} says]
\begin{enumerate}
  \item At fixed variance share, survival's precision share falls as $1/p^2$.
        Doubling the number of genes cuts it by four.
  \item It grows as $1/v_k^2$. Survival has the most relative influence on
        low-variance programs, which are the programs supervision is meant to
        find, and almost none on dominant programs.
  \item With the PDAC training data ($p\approx 2{,}200$, $D/n\approx 0.6$,
        $\betasec{k}\approx 0.04$, $\tau_j\approx 1.5$, $v_k\approx 0.05$) and
        $\ws = 1$, Eq.~\eqref{eq:rho-approx} gives $\rho\approx 10^{-6}$. The
        trace on the real data measured $10^{-8}$ to $10^{-6}$
        (DECISIONS.md, 2026-11-04). At that level the fitted loadings equal the
        unsupervised ones.
  \item Matching the balance across $p$ at fixed $v_k$ needs $\ws\propto p^2$.
        Fixed weights of 40 or more broke convergence at $p = 2{,}000$ in
        simulation, through separation of the Cox model. The comparison document
        (\texttt{SSBMF-LB-vs-YFB-comparison/}) discusses ways to give survival
        more weight without that failure.
\end{enumerate}
The approximation assumes signal-dominated scores and near-orthogonal programs.
For a program near the detection threshold the noise part of $Y\bar f_{\cdot k}$
is not negligible and Eq.~\eqref{eq:rho} should be used directly.
\end{keybox}

\paragraph{Comparison with the April 2026 claim.} The April derivation
concluded that the ratio is $O(1)$ because both precisions sum over $n$. It
treated $\sum_i\lsec{i}{k}$ as $n$ times a fixed variance, but the size of $l_k$
is tied to the size of $f_k$ through the reconstruction and the score scale.
Once that link is used, the ratio carries the factor $1/(p^2v_k^2)$.

% =============================================================================
\section{Strengthening Supervision Without Breaking the Fit}\label{sec:strengthen}
% =============================================================================

Section~\ref{sec:balance} shows that at $\ws = 1$ survival has almost no pull on
the loadings of a dominant program at $p\approx 2{,}000$. This section lists the
ways to change the balance and what each does to the derivation. None is
adopted.

\paragraph{Why a large weight fails.} The survival-only estimate in
Eq.~\eqref{eq:F-combination} solves, for program $k$, a weighted least squares
problem $\beta_k Y\bm f_{\cdot k}\approx\bm z^{-k}$ with $p$ unknowns and $n$
observations. With $p > n$ it has exact solutions: some combination of genes
orders the training survival times perfectly. As $\rho\to\infty$ the loadings
follow that solution, the Cox fit separates, $|\betabar{k}|$ grows without
bound and the program detaches from expression structure. This matches the
S-9 result that fixed weights of 40 or more lose convergence and the program.
Any remedy has to keep the genomics term or the priors in control of the
directions the Cox model alone cannot identify.

\begin{center}
\begin{tabularx}{\textwidth}{>{\raggedright\arraybackslash}p{3.3cm}>{\raggedright\arraybackslash}X>{\raggedright\arraybackslash}p{4.2cm}}
\toprule
Approach & Change to the derivation & Status / risk \\
\midrule
Weight $\ws$ on the Cox term (current) & $\ws$ in Eqs.~\eqref{eq:AB-beta} and \eqref{eq:AB-F}; power posterior & S-9: $\ws\ge 40$ separates at $p=2{,}000$ \\
Survival-block weight (App.~\ref{app:options}) & $\ws$ also multiplies $\KL(q_\beta\Vert g_\beta)$; $\beta$ update at weight 1 & Keeps the $\beta$ prior in force; CV choice stable (S-13), criteria not all met \\
Choose $\ws$ from data & Same updates; $\ws$ by cross-validated partial log-likelihood, or a calibrated power (generalized Bayesian updating) & CV picks 1--10 in simulation; the candidate for the next real-data refit \\
Divide the genomics term by $p$ & Equivalent to $\ws = p$ with the $L$, $F$ KL terms unscaled & $\rho\propto 1/p$, still falls with $p$; already inside the separating range at $p\approx 2{,}000$ \\
Normalize both losses and mix with $\alpha$ (DeSurv) & $(1-\alpha)\cdot$genomics$/c_G + \alpha\cdot$Cox$/c_S$ equals $\ws = \tfrac{\alpha}{1-\alpha}\tfrac{c_G}{c_S}$ up to an overall scale & Same family as $\ws$; DeSurv picks $\alpha$ by CV C-index \\
Fewer genes & $\rho\propto 1/p^2$, so a smaller gene set raises survival's share & A survival screen must sit inside CV to avoid reusing outcomes \\
Program-specific weights $w_{S,k}$ & $w_{S,k}$ in program $k$'s $F$ and $\beta$ updates & The Cox term couples programs, so this no longer maximizes one objective; heuristic only \\
\bottomrule
\end{tabularx}
\end{center}

\begin{keybox}[What the balance analysis implies for the method]
Because $\rho\propto 1/v_k^2$, a modest weight already acts on low-variance
programs and leaves dominant programs to the genomics term. That is the regime
where a joint fit can find what a two-step fit misses, and the near-threshold
simulation (S-14) is where supervision showed a gain (validation C $+0.156$,
5 of 10 recoveries against 0). Choosing $\ws$ by cross-validation targets that
regime without moving into separation. A larger $\ws$ cannot make survival
dominate a high-variance program without first breaking the fit.
\end{keybox}

% =============================================================================
\section{The Historical Configuration as a Special Case}\label{sec:historical}
% =============================================================================

The benchmark fits reported through October 2026 used
\texttt{alpha\_F = 0} with \texttt{canonicalize = FALSE}. The $F$ update then
mixed the two sources as $(1-\alpha_F)\cdot\text{genomics} + \alpha_F\cdot\text{survival}$,
and at $\alpha_F = 0$ the survival part has weight zero.

\begin{keybox}[At $\alpha_F = 0$ the fit is a two-step method]
With no survival term in the $F$ update, none of the $L$, $F$ or $\tau$ updates
uses $(t,\delta)$ or $\bm\beta$. The sequence of $(q_L, q_F, \tau)$ is the one an
unsupervised EBMF fit would produce from the same start. The $\beta$ update is
an empirical Bayes Cox regression on the current scores. The result is
unsupervised EBMF followed by a shrinkage Cox fit, which is why the benchmark
results did not differ from a two-step analysis. It remains a useful comparator.
\end{keybox}

\paragraph{How the historical weights differ from $\ws$.}
\begin{center}
\begin{tabularx}{\textwidth}{>{\raggedright\arraybackslash}p{3.1cm}>{\raggedright\arraybackslash}X>{\raggedright\arraybackslash}X}
\toprule
 & Historical ($\alpha$, $\alpha_F$) & Current ($\ws$) \\
\midrule
$F$ update & $(1-\alpha_F)A^{\text{gen}} + \alpha_F A^{\text{surv}}$; $\alpha_F > 0$ also down-weights genomics; $B^{\text{surv}}$ incomplete & $A^{\text{gen}} + \ws A^{\text{surv}}$; genomics keeps weight 1; complete $B^{\text{surv}}$ \\
$\beta$ update & $A_{\beta_k} = \alpha\sum_i W_{ii}\zbar{i}{k}^2\cdot p$ (\texttt{per\_p} divisor); no $V^Z$ & $A_{\beta_k} = \ws\sum_i W_{ii}\zsec{i}{k}$ \\
Score in the Cox model & $Y\bar f_{\cdot k}/\lVert\bar f_{\cdot k}\rVert$, while the $F$ update was derived for $Y\bar f_{\cdot k}$ & $Y\bar f_{\cdot k}$ with $\mathrm{sd} = 1$ \\
Monitored objective & genomics and survival weighted by $(1-\alpha)$ and $\alpha p$ & the same $\ws$ as the updates \\
Range & $\alpha,\alpha_F\in[0,1]$ & $\ws\in[0,\infty)$ \\
\bottomrule
\end{tabularx}
\end{center}
In both, the weight on $\beta$ cancels in $x_k$ and only changes $s_k$. The
historical code also flips the sign of $\bar{\bm\beta}$ after fitting if the
training C-index is below 0.5 (\texttt{sign\_correction = TRUE}, the default).
This is a post-processing heuristic, not part of the derivation, and the
supervision studies set it to \texttt{FALSE}.

% =============================================================================
\section{Extension to $M$ Modalities}\label{sec:multimodal}
% =============================================================================

The single-modality model is the case $M = 1$ of a model with shared patient
scores and modality-specific loadings (expression and methylation, for example):
\[
  Y_m = LF_m^\top + E_m,\quad E_{m,ij}\sim\mathcal N(0,\tau_{mj}^{-1}),\qquad
  \eta_i = \sum_k z_{ik}\beta_k,\quad
  Z = \sum_{m=1}^M Y_mF_m .
\]
The full derivation is in \texttt{SSBMF-YFB-multimodal-model-derivations/}. Every
update keeps the form derived here:
\begin{itemize}
  \item $q(l_k)$: $A_{l_k} = \sum_m\sum_j\tau_{mj}\overline{f^2_{m,jk}}$ and
        $B_{l_{ik}} = \sum_m\sum_j\tau_{mj}R^{-k}_{m,ij}\bar f_{m,jk}$. Still
        genomics only.
  \item $q(f_{m,k})$: Eq.~\eqref{eq:AB-F} with $Y_m$, $\tau_m$ and the joint score
        $\bar z_{\cdot k}$, which every modality's sweep refreshes. Each $F_m$ is
        dual-source.
  \item $q(\beta_k)$: Eq.~\eqref{eq:AB-beta} with the joint score and
        $V^Z_{ik} = \sum_m\sum_j y_{m,ij}^2\Var(f_{m,jk})$.
  \item $\tau_m$: Eq.~\eqref{eq:tau} for each modality.
  \item Score standardization: $\mathrm{sd}_i(\bar z_{ik}) = 1$ for the joint score.
\end{itemize}
With $M = 1$ every formula reduces to the one in this document.

\paragraph{New balance questions.} The precision of $l_k$ sums over all
$p = \sum_m p_m$ features, so the modality with more features (methylation)
dominates the shared scores. Survival's share of each loading precision falls
as in Eq.~\eqref{eq:rho-approx}, with the total $p$. Both point to
modality-specific weights on the reconstruction terms, which need their own
derivation and simulation check.

\paragraph{Convergence: what failed and what has changed since.}
\begin{center}
\begin{tabularx}{\textwidth}{>{\raggedright\arraybackslash}p{4.4cm}>{\raggedright\arraybackslash}X}
\toprule
Earlier multimodal problem & Single-modality change that addresses it \\
\midrule
Scale drift along $(a_kL_k, F_{mk}/a_k, a_k\beta_k)$; latent-moment stopping rules never fired (2026-09-18) & Score standardization each iteration (already in the multimodal fitter), and a scale-invariant objective \\
Stopping on relative change in reconstruction and predictor, with no recorded ELBO, so ELBO-based criteria could not be assessed (2026-10-15) & A monitored objective with the same weights as the updates and an ELBO stopping rule (Sections~\ref{sec:elbo}, \ref{sec:other}) \\
Loading KL entered with the wrong sign through the shared helper (2026-10-15) & Fixed, with a test that recomputes the last objective from the returned moments \\
Real-data fits reached the sweep cap; refitted priors could have a support incompatible with the cached posterior KL (2026-10-15) & Prior held fixed during each sweep and refit once after it; a dead-program rule replaces revived zero programs \\
Free noise precisions ran away as reconstruction error grew (S-8) & Inputs standardized within platform; S-9 canonical fits converged in every case. A fixed or bounded $\tau$ remains a fallback \\
Incomplete $B^{\text{surv}}$ in early single-modality code & The multimodal sweep already used the complete term; both fitters now share it \\
\bottomrule
\end{tabularx}
\end{center}

\paragraph{Recommended path.}
\begin{enumerate}
  \item Meet the single-modality acceptance criteria (recovery, false activity,
        convergence, external prediction) with this algorithm.
  \item Build one fitter for $M$ modalities whose $M = 1$ path reproduces
        \texttt{fit\_ssbmf\_yfb(canonicalize = TRUE)} to numerical tolerance.
        This is a direct test that the extension changes nothing for one block.
  \item Carry over the monitored objective with $\ws$, the ELBO stopping rule, the
        dead-program rule and the KL tests. Record the full ELBO at every
        iteration.
  \item Run two-modality simulations with planted programs in each modality and
        one shared prognostic program, then derive and test modality weights if
        methylation dominates $L$.
\end{enumerate}

% =============================================================================
\section{Differences from the April 2026 Derivations}\label{sec:errata}
% =============================================================================

\begin{center}
\begin{tabularx}{\textwidth}{c>{\raggedright\arraybackslash}p{3.4cm}>{\raggedright\arraybackslash}X>{\raggedright\arraybackslash}p{2.4cm}}
\toprule
ID & Location (\texttt{components/}) & Issue and correction & Here \\
\midrule
Y1 & \texttt{qF\_supervised}, eq.\ BFsurv & $B^{\text{surv}}$ removed the other programs but not the other genes of program $k$. Correct term subtracts $\betasec{k}\sum_i W_{ii}y_{ij}(\zbar{i}{k}-y_{ij}\fbar{j}{k})$, and the update becomes a sequential sweep (found 2026-10-15). & Eq.~\eqref{eq:AB-F-surv} \\
Y2 & \texttt{qF\_supervised}, eqs.\ AF, BF & Convex mixing $(1-\alpha)/\alpha$ replaced by weight $\ws$ on the Cox term; genomics weight 1. & Eq.~\eqref{eq:AB-F} \\
Y3 & \texttt{qF\_supervised}, Sec.\ 4 & ``Scale imbalance resolved, ratio $O(1)$'' is incorrect; the ratio scales as $1/(p^2v_k^2)$. & Sec.~\ref{sec:balance} \\
Y4 & \texttt{qF\_supervised}, Sec.\ 5 & ``EBNM shrinkage identifies the scale; no constraint needed.'' The free scale fed back into the survival terms; fixed by score standardization. & Sec.~\ref{sec:scaling} \\
Y5 & \texttt{qBeta}, eqs.\ Abeta & Scores treated as observed: $A_{\beta_k}$ used $\zbar{i}{k}^2$. It should use $\zsec{i}{k} = \zbar{i}{k}^2 + V^Z_{ik}$. & Eq.~\eqref{eq:AB-beta} \\
Y6 & \texttt{qBeta}, Sec.\ 2 & Cox model used $Y\bar F D^{-1}$ (unit-norm columns), off the scale the $F$ update was derived for. & Sec.~\ref{sec:scaling} \\
Y7 & \texttt{ELBO}, term (II) & $\Var(\eta_i)$ omitted the loading-uncertainty part $V^Z_{ik}\betasec{k}$. & Eq.~\eqref{eq:surv-expect} \\
Y8 & \texttt{ELBO} / code & Monitored objective weighted survival differently from the updates, and the loading KL entered with the wrong sign (fixed 2026-10-15). & Sec.~\ref{sec:other} \\
\bottomrule
\end{tabularx}
\end{center}

The $q(L)$ derivation (\texttt{qL\_unsupervised}) and the $\tau$ update are
unchanged.

\appendix
% =============================================================================
\section{Options Tested and Not Adopted}\label{app:options}
% =============================================================================

\paragraph{Survival-block weight (\texttt{survival\_block\_weight = TRUE}).}
$\ws$ multiplies $\KL(q_\beta\Vert g_\beta)$ as well as the Cox term:
\[
  \mathcal F_1 = \E[\log p(Y\mid L,F,\tau)] - \KL_L - \KL_F
  + \ws\bigl[\E\,\ell_{\text{Cox}} - \KL(q_\beta\Vert g_\beta)\bigr].
\]
The $L$, $F$ and $\tau$ updates are those of Eq.~\eqref{eq:elbo}. For $\ws > 0$
the maximizer over $(q_\beta, g_\beta)$ no longer depends on $\ws$, and the
$\beta$ update uses weight 1 in Eq.~\eqref{eq:AB-beta}. This keeps the prior on
$\beta$ at full weight however large $\ws$ is. In simulation (S-13), choosing
$\ws$ by cross-validated partial log-likelihood under this objective was stable
but did not pass all registered criteria.

\paragraph{Joint Gaussian $q(\bm\beta)$ (\texttt{beta\_posterior = "joint"}).}
Replace $\prod_k q(\beta_k)$ by $\mathcal N_K(\bm m,\Sigma)$ with
$G = \ws\bigl(\bar Z^\top W\bar Z + \mathrm{diag}(\sum_i W_{ii}V^Z_{ik})\bigr)$,
$\bm b = \ws\bar Z^\top W\bm z$, $\Sigma = (G+\Lambda)^{-1}$, $\bm m = \Sigma\bm b$,
and per-program prior variances by type-II maximum likelihood. For $K = 1$ this
equals the normal-family EBNM update. In S-13 it gave the same results as the
mean-field update.

\paragraph{Shared prior on $\bm\beta$ (\texttt{beta\_prior\_mode = "shared"}).}
One prior $g_\beta$ for all programs, held fixed during the sweep and refit from
the $K$ pseudo-observations after it, in place of $K$ priors each fit to one
observation.

\end{document}
